% !TEX root = ../main.tex
\corrigesChapitre{14}{Automatismes et synthèses}
\corrigesSouples

\begin{blocLivre}{cor-14-1}
\begin{corrige}{1}{Évolutions et retour au départ}
\demarche{Composer les coefficients multiplicateurs plutôt que les taux.}
\begin{questions}
\item $80\times1{,}25=100$ euros, puis $100\times0{,}8=\res{80}$ euros. Le coefficient global vaut 1, donc le taux global est $\res{0\,\%}$.
\item Le coefficient réciproque est $1/0{,}8=1{,}25$, soit une hausse de $\res{25\,\%}$.
\end{questions}
\end{corrige}
\end{blocLivre}
\begin{blocLivre}{cor-14-2}
\begin{corrige}{2}{Retrouver une valeur et un taux}
\demarche{Retrouver la valeur initiale par division et calculer les taux relativement à cette valeur.}
\begin{questions}
\item $0{,}7P=84$, donc $\res{P=120\text{ euros}}$.
\item $(54-45)/45=9/45=\res{20\,\%}$.
\item $1{,}1\times1{,}2=1{,}32$, donc le taux global est $\res{32\,\%}$.
\end{questions}
\end{corrige}
\end{blocLivre}
\begin{blocLivre}{cor-14-3}
\begin{corrige}{3}{Développer, factoriser, résoudre}
\demarche{Reconnaître l’identité remarquable ou le produit nul adapté.}
\begin{questions}
\item $\res{4x^2-12x+9}$.
\item Différence de carrés : $\res{(3x-5)(3x+5)}$.
\item Un facteur doit être nul : $\res{x=2\text{ ou }x=-1}$.
\end{questions}
\end{corrige}
\end{blocLivre}
\begin{blocLivre}{cor-14-4}
\begin{corrige}{4}{Signe et intervalles}
\demarche{Distinguer le signe d’un facteur et celui d’un produit.}
\begin{questions}
\item On obtient $\res{x\leq3/2}$ en divisant par le coefficient négatif $-2$ et en inversant le sens de l’inégalité.
\item Le produit est négatif strictement entre les racines : $\res{x\in\mathopen{]}-2;3[}$.
\item Par exemple $\res{x=0}$ satisfait les deux conditions. L’ensemble commun est $\mathopen{]}-2;3/2]$.
\end{questions}
\end{corrige}
\end{blocLivre}
\begin{blocLivre}{cor-14-5}
\begin{corrige}{5}{Équations de droites}
\demarche{Calculer la pente avec deux points, puis utiliser l’un d’eux pour l’ordonnée à l’origine.}
\begin{questions}
\item La pente vaut $(7-3)/(3-1)=2$. Avec A, $3=2+p$, donc $\res{y=2x+1}$.
\item La pente est $-1/2$, l’ordonnée à l’origine 4. $y=0$ donne $x=8$, donc l’intersection est $\res{(8;0)}$.
\end{questions}
\label{figlivre-r14-05-droites-cor}
\figureLivre[corrige]{r14-05-droites}

\end{corrige}
\end{blocLivre}
\begin{blocLivre}{cor-14-6}
\begin{corrige}{6}{Lire une courbe par morceaux}
\demarche{Repérer les intersections avec la droite horizontale et les portions au-dessus.}
\begin{questions}
\item Les intersections avec $y=1$ sont $\res{x=-2;1;4}$. L’inéquation donne $\res{\{-2\}\cup[1;4]}$. La valeur 1 sur le segment croissant se vérifie avec $y=2x-1$.
\item Décroissance sur $[-2;0]$, croissance sur $[0;2]$, décroissance sur $[2;4]$. Minimum $-1$ en 0, maximum 3 en 2.
\end{questions}
\label{figlivre-r14-06-lecture-cor}
\figureLivre[corrige]{r14-06-lecture}

\end{corrige}
\end{blocLivre}
\begin{blocLivre}{cor-14-7}
\begin{corrige}{7}{Des fréquences conditionnelles}
\demarche{La population de référence change entre fréquence globale et fréquence conditionnelle.}
\begin{questions}
\item Il y a $20+30=50$ utilisateurs : fréquence $\res{1/2}$.
\item Parmi les 50 plus jeunes, la fréquence est $20/50=\res{2/5}$.
\item La probabilité conditionnelle $2/5$ diffère de la probabilité globale $1/2$ : pas d’indépendance dans ce modèle.
\end{questions}
\label{figlivre-r14-07-cor}
\figureLivre[corrige]{r14-07}

\end{corrige}
\end{blocLivre}
\begin{blocLivre}{cor-14-8}
\begin{corrige}{8}{Réactiver la trigonométrie}
\demarche{Retirer un tour pour comparer les points-images, et utiliser les radians pour les arcs.}
\begin{questions}
\item $5\pi/3=2\pi-\pi/3$ : $\res{(1/2;-\sqrt3/2)}$.
\item La différence entre les deux angles vaut $2\pi$ : les coordonnées sont les mêmes.
\item $L=R\theta=3\times2\pi/3=\res{2\pi}$.
\end{questions}
\label{figlivre-r14-08-trigo-cor}
\figureLivre[corrige]{r14-08-trigo}

\end{corrige}
\end{blocLivre}
\begin{blocLivre}{cor-14-9}
\begin{corrige}{9}{Relier dérivée et extremums}
\demarche{Factoriser la dérivée et comparer tous les candidats, y compris les bornes.}
\begin{questions}
\item $f\prime(x)=3x^2-12x+9=\res{3(x-1)(x-3)}$, de signes $+$, $-$, $+$ successivement.
\item $f$ croît de $f(0)=0$ à $f(1)=4$, décroît à $f(3)=0$, puis croît à $f(4)=4$.
\item Minimum $\res{0}$ en 0 et 3 ; maximum $\res{4}$ en 1 et 4. Il faut comparer aussi les bornes.
\end{questions}
\label{figlivre-r14-09-var-cor}
\figureLivre[corrige]{r14-09-var}

\end{corrige}
\end{blocLivre}
\begin{blocLivre}{cor-14-10}
\begin{corrige}{10}{Un diagramme circulaire}
\demarche{Convertir les fréquences en angles tout en conservant le total de 360 degrés.}
\begin{questions}
\item Les fréquences sont $\res{1/2;3/10;1/5}$. Les angles sont $\res{180^\circ;108^\circ;72^\circ}$ et totalisent $360^\circ$.
\item $40-24=\res{16\text{ voix}}$.
\item Non : les noms sont des catégories qualitatives, sans valeur numérique à moyenner. On peut comparer leurs effectifs ou leurs fréquences.
\end{questions}
\label{figlivre-r14-10-secteurs-cor}
\figureLivre[corrige]{r14-10-secteurs}

\end{corrige}
\end{blocLivre}
\begin{blocLivre}{cor-14-11}
\begin{corrige}{11}{Une exponentielle dans une étude de fonction}
\demarche{Factoriser la dérivée par une exponentielle positive.}
\begin{questions}
\item Par injectivité, $2x=x+3$, donc $\res{x=3}$.
\item $g\prime(x)=2e^{-x}-(2x-1)e^{-x}=\res{(3-2x)e^{-x}}$.
\item La fonction croît jusqu’à $3/2$, puis décroît. Le maximum vaut $\res{2e^{-3/2}}$. Aux bornes, $g(0)=-1$ et $g(3)=5e^{-3}>0$ : minimum $\res{-1}$ en 0.
\end{questions}
\label{figlivre-r14-11-var-cor}
\figureLivre[corrige]{r14-11-var}

\end{corrige}
\end{blocLivre}
\begin{blocLivre}{cor-14-12}
\begin{corrige}{12}{Synthèse : une boîte à dimensions liées}
\demarche{Relier le schéma du solide à son volume, puis étudier un signe sur le domaine géométrique.}
\begin{questions}
\item La base carrée a pour aire $x^2$ et la hauteur vaut $6-x$.
Donc $V(x)=x^2(6-x)=-x^3+6x^2$ et $V'(x)=-3x^2+12x=\res{3x(4-x)}$.
\item Sur $]0;6[$, le facteur $3x$ est positif. La dérivée a le signe de $4-x$ :
le maximum est atteint en $x=4$, avec $V(4)=16\times2=\res{32\text{ L}}$.
La base mesure $4\times4$ dm et la hauteur 2 dm ; $1$ dm$^3=1$ litre.
\item $V(2)=16$ et $V'(2)=12$. L’équation de la tangente est
$y=12(x-2)+16=\res{12x-8}$.
\item En développant, $-(x-2)(x^2-4x-8)=-x^3+6x^2-16=V(x)-16$.
Le second facteur a pour racines $2-2\sqrt3<0$ et $2+2\sqrt3\in]0;6[$.
Sur le domaine géométrique, le tableau de signes ci-dessous donne
$V(x)-16\geq0\iff\res{x\in[2;2+2\sqrt3]}$.
Les deux bornes sont incluses, car elles donnent exactement 16 litres.
\end{questions}

\label{figlivre-r14-12-var-cor}
\figureLivre[corrige]{r14-12-var}

\label{figlivre-r14-12-signes-cor}
\figureLivre[corrige]{r14-12-signes}

\label{figlivre-r14-12-boite-cor}
\figureLivre[corrige]{r14-12-boite}

\end{corrige}
\end{blocLivre}
\begin{blocLivre}{cor-14-13}
\begin{corrige}{13}{Synthèse : triangle, cercle et tangente}
\demarche{Relier produit scalaire, cercle de diamètre donné, tangente et longueur d’arc.}
\begin{questions}
\item $\vect{CA}=(-2;-2)$ et $\vect{CB}=(2;-2)$ ont pour produit $-4+4=0$ : l’angle en $C$ est droit.
De plus $AC^2=BC^2=8$, donc le triangle est rectangle isocèle.
Ses deux autres angles sont égaux et se partagent $\pi/2$ : $\res{\widehat{BAC}=\pi/4}$.
\item Dans un triangle rectangle, le milieu de l’hypoténuse est le centre du cercle circonscrit.
Ici le diamètre est $[AB]$, le centre $\Omega(2;0)$ et le rayon 2 :
$\res{(x-2)^2+y^2=4}$.
\item Le rayon $\Omega C$ est vertical, donc la tangente en $C$ est horizontale : $\res{y=2}$.
La projection verticale de $M(5;3)$ conserve son abscisse et prend l’ordonnée de la tangente : $\res{H(5;2)}$.
\item Les rayons $\vect{\Omega A}=(-2;0)$ et $\vect{\Omega C}=(0;2)$ sont perpendiculaires.
L’angle central du petit arc vaut donc $\pi/2$ radians.
Sa longueur est $R\theta=2(\pi/2)=\res{\pi}$ unités de longueur.
\end{questions}

\label{figlivre-r14-13-geometrie-cor}
\figureLivre[corrige]{r14-13-geometrie}

\end{corrige}
\end{blocLivre}
\begin{blocLivre}{cor-14-14}
\begin{corrige}{14}{Synthèse : défauts et gain net}
\demarche{Lier l’arbre des fournisseurs à la loi du gain net.}
\begin{questions}
\item Les deux fournisseurs forment une partition.
$P(A\cap D)=0{,}6\times0{,}01=0{,}006$ et $P(B\cap D)=0{,}4\times0{,}04=0{,}016$.
Donc $P(D)=0{,}006+0{,}016=\res{0{,}022}$.
\item Parmi les articles défectueux, la part provenant de B vaut
$P_D(B)=P(B\cap D)/P(D)=0{,}016/0{,}022=\res{8/11}$.
\item Un article sans défaut donne $20-12=8$ euros ; avec défaut, il donne $20-12-30=-22$ euros.
Le tableau récapitule cette loi.
Ainsi $E(G)=8(0{,}978)-22(0{,}022)=\res{7{,}34\text{ euros}}$.
Le deuxième moment vaut $E(G^2)=64(0{,}978)+484(0{,}022)=73{,}24$,
d’où $V(G)=73{,}24-7{,}34^2=\res{19{,}3644}$ euros carrés.
\item L’évènement contraire est aucun défaut sur les quatre articles.
Par indépendance, il a pour probabilité $0{,}978^4$.
On obtient $P(\text{au moins un défaut})=\res{1-0{,}978^4\approx0{,}08514}$.
\end{questions}

\label{figlivre-r14-14-cor}
\figureLivre[corrige]{r14-14}

\label{figlivre-r14-14-loi-cor}
\figureLivre[corrige]{r14-14-loi}

\end{corrige}
\end{blocLivre}
\begin{blocLivre}{cor-14-15}
\begin{corrige}{15}{Synthèse : deux modèles d'évolution}
\demarche{Comparer les deux modèles aux instants entiers en vérifiant les seuils et leurs prédécesseurs.}
\begin{questions}
\item $v_{n+1}=0{,}9v_n$ et $v_0=500$, donc $\res{u_n=500+500(0{,}9)^n}$.
\item $u_{n+1}-u_n=-50(0{,}9)^n<0$ et $u_n-500=500(0{,}9)^n>0$.
\item $u_{15}\approx602{,}946>600$ et $u_{16}\approx592{,}651\leq600$ : premier rang $\res{16}$.
\item $W\prime(t)=-50e^{-0{,}1t}<0$. $W(16)\approx600{,}948$ et $W(17)\approx591{,}342$ : première semaine entière $\res{17}$.
\item Pour $x=-0{,}1\neq0$, $e^{-0{,}1}>0{,}9$. En élevant à la puissance entière positive n, $(e^{-0{,}1})^n>(0{,}9)^n$, donc $\res{W(n)>u_n}$. Le modèle continu décroît moins vite aux instants entiers, d’où le seuil plus tardif.
\end{questions}

Pour justifier que les seuils trouvés sont les premiers, on utilise la décroissance de chaque modèle, puis les valeurs de part et d’autre : $u_{15}>600\geq u_{16}$ et $W(16)>600\geq W(17)$.

\label{figlivre-r14-15-modeles-cor}
\figureLivre[corrige]{r14-15-modeles}

\end{corrige}
\end{blocLivre}
